\documentclass[12pt]{article}

\usepackage[paper=letterpaper,margin=2.5cm]{geometry} % Set Margins

%% Math and math fonts
\usepackage{amsmath, amsthm, amssymb, amsfonts}
\usepackage{bbold} % for \mathbbm{1}

% date
\usepackage[nodayofweek]{datetime}

% Color
\usepackage{color, xcolor}

% Misc
\usepackage{environ}  % \collect@body in asmmath
\usepackage{graphicx} % \includegraphics options
\usepackage{mdframed} % text boxes
\usepackage{indentfirst} % Indent first paragraph after section header
\usepackage[shortlabels]{enumitem} % Control enumerate items with [(a)]
\usepackage{comment} % Comments
\usepackage{fancyhdr} % Headers and footers

% Tables
\usepackage{array}

% Sub-figures and figure placement
\usepackage{caption}
\usepackage{subcaption}
\usepackage{float} 

% Graphing
\usepackage{pgfplots}
\pgfplotsset{compat=1.17}
\usepackage{tikz}

% Title Placement
\usepackage{titling}
\setlength{\droptitle}{-6em}

%set indent to 
\setlength{\parindent}{0pt}
\setlength{\parskip}{10pt}

% Hyper refs
\usepackage{hyperref}
\hypersetup{
    colorlinks=true,
    linkcolor=blue,
    urlcolor  = blue,
    filecolor=magenta,      
    urlcolor=blue,
    citecolor = blue,
    anchorcolor = blue
}

% % Citation management
\usepackage{natbib}
\bibliographystyle{abbrvnat}
\setcitestyle{authordate,open={(},close={)}}

% ----------------------------------------
% TITLE
% ----------------------------------------

\pagestyle{fancy}

\lhead{Stone}
\chead{Class One}
\rhead{FAME}

\title{FAME Class 1 Notes -- Mathematics as Creative Symbolic Reasoning}
\author{Eliana Stone}

\begin{document}
\maketitle
\section{What is the goal of this course?}
This course aims to give you the tools to formalize problems and ``speak math'', so that you can understand and engage more fully with interdisciplinary problems elsewhere. More practically, this course will also provide the math background needed to access a wider range of courses at Yale that require some math familiarity. 

This class will help you choose which ``math'' approach should be applied to a given problem. We don't require you to memorize equations for regurgitation. Instead, you are familiarized with solution strategies and trained in how to apply them to environmental problems as they come up in various facets of your life and work.

\section{Canvas Overview and Course Objectives}

See: \hyperlink{Canvas}{https://yale.instructure.com/courses/121587}

\section{Introduction: Science vs Math}
Science is based on empirical evidence and observation (data). Math is based on logical reasoning and proof. In science, we collect data through repeated observations and experiments to test hypotheses, but we cannot prove/conclude anything. In math, you do not need much evidence. Instead, we begin with a set of assumptions (often simple and abstract) and use logical reasoning to prove conclusions that follow from those assumptions. Essentially, science uses observations to test hypotheses and rule out alternative explanations, and math uses proofs to establish conclusions within an internally consistent framework.

To establish conclusions, math relies on understanding relationships between numbers, variables, functions, and other abstract objects. To do this, we need notation to express these relationships clearly. Math is also about communicating, so not only do you as an individual need mathematical notation that makes sense to you, but also to whomever you are communicating with.

\section{Numbers}

Numbers can be thought of as adjectives. They describe quantities, characteristics, or relationships between objects.

For example, we might use numbers to describe the number of trees in a forest, the temperature of a lake, or the amount of rainfall in a particular location.

Different math problems require different types of numbers. These numbers are organized into sets. We'll discuss some useful sets of numbers to know.

\subsection{Real Numbers}

Real numbers are denoted by $\mathbb{R}$ and include all numbers along the number line:

\[
\mathbb{R}=(-\infty,\infty)
\]

This includes whole numbers, fractions, decimals, negative numbers, and irrational numbers.

For example,

\[
-3,\quad 0,\quad \frac{1}{2},\quad \sqrt{2},\quad \pi
\]

are all real numbers.

Real numbers are continuous, meaning that between any two distinct real numbers there are infinitely many additional real numbers.

For example, between 0 and 1, we have

\[
0.1,\quad 0.01,\quad 0.001,\quad 0.5,\quad \ldots
\]

Importantly, infinity ($\infty$) and negative infinity ($-\infty$) are not themselves real numbers. They describe unbounded directions along the number line.

\subsection{Positive Real Numbers}

It is often useful to restrict the real numbers to particular ranges.

Positive real numbers include zero and every real number greater than zero:

\[
\mathbb{R}^+=[0,\infty)
\]

You might also see a set of strictly positive real numbers that include every real number greater than zero, but exclude zero itself:

\[
\mathbb{R}^+=(0,\infty)
\]

Notice the distinction between brackets and parentheses:

\begin{itemize}
    \item Square brackets $[\,]$ indicate that an endpoint is included.
    \item Parentheses $(\,)$ indicate that an endpoint is excluded.
\end{itemize}

\subsection{Real Numbers in Multiple Dimensions}

So far, we have discussed real numbers along a single number line. However, math problems often involve multiple dimensions at once.

For example, a point in two-dimensional space can be described by two coordinates:

\[
(x,y)
\]

The set of all ordered pairs of real numbers is denoted by

\[
\mathbb{R}^{++}=\mathbb{R}^+\times\mathbb{R}^+
\]

If both coordinates must be nonnegative, we write

\[
\mathbb{R}_{\geq 0}^2
=
[0,\infty)\times[0,\infty)
\]

This means

\[
x\geq 0
\qquad \text{and} \qquad
y\geq 0
\]

Similarly, $\mathbb{R}^3$ represents three-dimensional space, consisting of ordered triples $(x,y,z)$.

In general, $\mathbb{R}^n$ describes an $n$-dimensional space.

\subsection{Natural Numbers and Integers}

Natural numbers are our counting numbers:

\[
\mathbb{N}=\{1,2,3,4,\ldots\}
\]

Some mathematical conventions include zero in the natural numbers. In this course, unless otherwise specified, we will use the convention that natural numbers begin at 1.

Unlike real numbers, natural numbers are discrete. There is no natural number between 1 and 2.

Integers include all positive and negative whole numbers, as well as zero:

\[
\mathbb{Z}=\{\ldots,-3,-2,-1,0,1,2,3,\ldots\}
\]

Natural numbers are a subset of integers, which are themselves a subset of real numbers:

\[
\mathbb{N}\subseteq\mathbb{Z}\subseteq\mathbb{R}
\]

The symbol $\subseteq$ means ``is a subset of.''

\section{Units}

\textbf{PAY ATTENTION TO UNITS!}

Math allows us to describe relationships between quantities, but those quantities often represent real-world measurements.

When working with interdisciplinary teams, it is especially important to make sure everyone is using the same units. Sometimes, apparent disagreements are simply the result of people expressing the same quantity in different units.

For example, a distance of 1 kilometer is equivalent to 1,000 meters:

\[
1\text{ km}=1{,}000\text{ m}
\]

These numbers are different, but still represent the same physical distance.

Units also help us determine whether an equation makes sense. For example, we cannot meaningfully add a distance measured in meters to a time measured in seconds. This is one of the simplest ways to identify mistakes in mathematical and scientific calculations.


\section{Functions}

\subsection{Functions Map Inputs to a Single Output}

A function is a mathematical relationship that assigns exactly one output to each input.

We can write

\[
y=f(x)
\]

Here:

\begin{itemize}
    \item $x$ is the input.
    \item $f$ is the function, or the rule describing the relationship.
    \item $f(x)$ is the output associated with input $x$.
    \item $y$ is equal to that output.
\end{itemize}

Importantly, a function does not necessarily describe a process that creates/builds something. Instead, it specifies a relationship between inputs (like $x$ above) and outputs (like $y$ above), rather than describing an action that constructs $y$. Additionally, functions are mathematical representations of relationships. Whether these representations accurately describe relationships in the real world is an empirical question that requires scientific evidence.

\subsection{A Note on Equal Signs}

The equal sign ($=$) means that the two things on either side of the equal sign will evaluate to the same thing (not that one constructs the other). With functions, we use the equal sign to say that $y$ \text{\underline{is}} $f(x)$.

For example without needing to think about functions,

\[
3+2=5
\]

We can also reverse the equation, so that

\[
5=3+2
\]

This is because equality is symmetric (or reversible), so if $a=b$, then $b=a$. 

Although equality works in both directions, this does not mean that when we say $y = f(x)$ every function can be uniquely reversed.

For example,

\[
f(x)=x^2
\]

Both $x=2$ and $x=-2$ produce the output 4. Knowing the output does not necessarily allow us to recover a unique input.

\subsection{The Vertical Line Test}

To determine whether a graph represents a function of $x$, imagine drawing a vertical line anywhere along the graph. If any vertical line intersects the graph more than once, the graph does not represent a function of $x$. This is because a function must assign exactly one output ($y$-value) to each input ($x$-value).

A function may assign the same output to multiple different inputs. What it cannot do is assign multiple outputs to the same input.

For example,

\[
y=x^2
\]

is a function because every value of $x$ produces exactly one value of $y$.

However, the equation

\[
x^2+y^2=1
\]

describes a circle and does not define $y$ as a function of $x$ over the entire circle.
For instance, when $x=0$, \(y=\pm 1\), so 
the same input produces two possible outputs, violating the definition of a function.

\subsection{Set Notation}

Before discussing functions more formally, we need to understand sets.

A \textbf{set} is a collection of distinct objects, called \textbf{elements}.

If $x$ is an element of set $X$, we write

\[
x\in X
\]

The symbol $\in$ means ``is an element of'' or ``belongs to.''

For example, suppose

\[
X=\{1,2,3,4\}
\]

Then, \(
2\in X
\)
but \(
5\notin X
\), where the symbol $\notin$ means "is not an element of" or "does not belong to".

We can also use subscripts to distinguish multiple elements within a set:

\[
x_1,x_2,x_3\in X
\]

Sets also do not have to contain numbers. They can contain animals, locations, categories, or other objects.

For example,

\[
P=\{\text{dog},\text{cat},\text{bird}\}, 
\text{ so dog}\in P. 
\]

We can also define sets that contain other sets. \\
\textbf{Note}: this is the more formal way of describing subsets as mentioned briefly earlier in these notes.

For example, if $D$ is the set of dogs and $A$ is the set of all animals, then

\[
D\subseteq A
\]

This means that every element of $D$ is also an element of $A$. 

Notice the distinction:
\begin{itemize}
    \item $\in$ describes membership, saying an element belongs to a set.
    \item $\subseteq$ describes inclusion, saying one set is contained within another.
\end{itemize}

\subsection{Domain, Codomain, and Range}

A function assigns exactly one output to every input in a specified set.

We can represent a function more formally as:

\[
f:X\rightarrow Y
\]

This reads: ``$f$ maps $X$ to $Y$.''

Here, $X$ is the domain, and $Y$ is the codomain.

\textbf{Domain:} The set of all permissible inputs to a function.

\textbf{Codomain:} The set in which the outputs of the function are defined to lie.

\textbf{Range:} The set of outputs that the function actually produces when evaluated over its entire domain.

The range is always a subset of the codomain, but it does not necessarily equal the codomain.

\subsection{Linearity}

One of the simplest and most useful mathematical relationships is a straight line.

The general equation of a straight line is

\[
\boxed{y=mx+b}
\]

where

\begin{itemize}
    \item $x$ is the input variable.
    \item $y$ is the output variable.
    \item $m$ is the slope of the line.
    \item $b$ is the $y$-intercept.
\end{itemize}

\textbf{Slope ($m$):} Describes how much $y$ changes when $x$ increases by one unit.

Mathematically,

\[
m=\frac{\Delta y}{\Delta x}
\]

where $\Delta$ means ``change in.''

The slope describes the rate of change of the output with respect to the input.

For a straight line, this rate of change is constant.

For example, consider

\[
y=3x+2
\]

If $x$ increases by one unit, $y$ increases by three units, regardless of the starting value of $x$.

We can see this by evaluating the function:

\begin{align*}
f(1)&=3(1)+2=5\\
f(2)&=3(2)+2=8\\
f(3)&=3(3)+2=11
\end{align*}

Each one-unit increase in $x$ produces a three-unit increase in $y$.

\textbf{Intercept ($b$):} Describes the value of $y$ when $x=0$.

For our example,

\[
f(0)=3(0)+2=2
\]

Therefore, the line crosses the $y$-axis at 2.

Linear models are a useful example of how math allows us to make up assumptions to simplify complicated problems. In the real world, relationships between variables are often messy and nonlinear. However, we do not need to necessarily know whether the relationship is linear in the real world in order to make the choice to represent it a linear way. That assumption can make the problem more manageable, easier to analyze, and more straightforward to communicate.

\subsection{Independent vs. Dependent variables}

When we write a function,

\[
y=f(x)
\]

we often describe $x$ as the independent variable and $y$ as the dependent variable.

In a purely mathematical setting, the independent variable serves as the input to the function. We specify its value and use the function to determine the corresponding output, and the dependent variable is mechanically determined by that corresponding output value. 

Colloquially however, we often use the words "independent" and "dependent" to imply something about causation. When we say that one variable is independent, we might imagine that it changes freely. Similarly, when we say that another variable is dependent, we might assume that its value changes \textit{because of} the independent variable. When we assign a scientific interpretation to the mathematical relationship expressed by the equal sign in $y=f(x)$, we need to bring in empirical evidence and additional assumptions beyond the math itself to make claims about causation.


\end{document}